\documentclass[12pt]{article}
\usepackage[margin=1in]{geometry}
\usepackage[T1]{fontenc}
\usepackage[utf8]{inputenc}
\usepackage{amsmath,amssymb,amsthm,mathtools,enumitem,mathrsfs}

\theoremstyle{plain}
\newtheorem{theorem}{Theorem}[section]
\newtheorem{lemma}[theorem]{Lemma}
\newtheorem{proposition}[theorem]{Proposition}
\newtheorem{corollary}[theorem]{Corollary}

\theoremstyle{definition}
\newtheorem{remark}[theorem]{Remark}

\numberwithin{equation}{section}

% Ground rings and fields
\newcommand{\Ring}{\mathbb{Z}[h^{\pm 1}]}
\newcommand{\Field}{\mathbb{Q}(h)}

% Skein algebra notation
\newcommand{\bigon}{\mathfrak{B}^{*}}
\newcommand{\hS}{\overline{\mathscr{S}}}

% Double algebra
\newcommand{\DqB}{D_qB}
\newcommand{\Khalf}{K^{1/2}}
\newcommand{\Kmh}{K^{-1/2}}
\newcommand{\Lhalf}{L^{1/2}}
\newcommand{\Lmh}{L^{-1/2}}
\newcommand{\Ebar}{\overline{E}}
\newcommand{\Fbar}{\overline{F}}
\newcommand{\bipr}{\bowtie}         % for the double cross product (Majid)
\newcommand{\Uqsc}{U_q^{\mathrm{sc}}(\mathfrak{sl}_2)}

\begin{document}

\title{The rank-two stated-skein realization of the generalized quantum double of the Borel subalgebras of \(U_q(\mathfrak{sl}_2)\)}
\author{}
\date{}
\maketitle

\begin{abstract}
We prove the \(n=2\) case of the stated-skein/quantum-double problem for the once-punctured bigon.  We construct an explicit algebra isomorphism from the rank-two generalized Drinfeld--Tanisaki double of the two quantum Borel subalgebras to the reduced stated-skein algebra of the once-punctured bigon.  In the arc dictionary, the lower \(b\)-arcs carry the positive Borel half and the upper \(c\)-arcs carry the negative Borel half.  We then show that quotienting the rank-two reduced stated-skein algebra by the ideal generated by the elements \(b_{ii}c_{ii}-1\) gives the simply connected rank-two quantum group, after the standard scalar extension needed to normalize Chevalley generators.  We also prove a torus-centrality criterion for doubles of this type, stated with the generation hypothesis needed for the proof.  No general-\(n\) skein isomorphism is asserted; the missing general-\(n\) ingredients are listed in the final open-points section.
\end{abstract}

\section{Introduction and statement of the rank‑two result}
\label{sec:intro}
% RECORDS: 41, 397, 699

The original problem asks for the following.

\begin{enumerate}[label=(\arabic*)]
\item Construct an explicit algebra isomorphism between the reduced stated-skein algebra \(\hS_n(\mathfrak B^*)\) of the once-punctured bigon and the generalized quantum double
\[
\check U_q(\mathfrak b_-)\,\bipr\,\check U_q(\mathfrak b_+)
\]
of the upper and lower Borel subalgebras of the simply connected quantum group \(\check U_q(\mathfrak{sl}_n)\), with respect to the standard Drinfeld--Tanisaki skew pairing
\[
\langle \widehat K_i,\widehat K_j\rangle
=
q^{\frac1n-\delta_{ij}},
\qquad
\langle E_i,F_j\rangle
=
\delta_{ij}/(q^{-1}-q),
\qquad
\langle \widehat K_i,F_j\rangle
=
\langle E_j,\widehat K_i\rangle
=
0.
\]
\item Deduce that \(\check U_q(\mathfrak{sl}_n)\) is isomorphic to
\[
\hS_n(\mathfrak B^*)/I,
\]
where \(I\) is the two-sided ideal generated by the elements \(b_{ii}c_{ii}-1\), with \(b_{ij}\) the lower rightward arcs and \(c_{kl}\) the upper rightward arcs.
\end{enumerate}
The results proved below establish the complete explicit statement only for \(n=2\).  Accordingly, this paper proves precisely the \(n=2\) specialization.  By ``rank-two'' we mean the problem's \(n=2\) case, although the Lie type is \(A_1\).

To avoid a common notational trap, we write \(h\) for the square-root parameter and put
\[
q=h^2.
\]
Korinman's parameter \(A\) in \cite{Korinman} is our \(h\).  We also use \(x,y,u,v\) for the surviving arc generators, so that no conflict arises with Korinman's parameter.
The main results are Theorem~\ref{thm:rank-two-double-skein} and Theorem~\ref{thm:rank-two-quotient}.  The first gives the explicit isomorphism between the rank-two generalized double and the reduced stated-skein algebra.  The second identifies the quotient by the rank-two ideal generated by \(b_{11}c_{11}-1\) and \(b_{22}c_{22}-1\) with the simply connected quantum group \(\Uqsc\), after the scalar extension needed to normalize Chevalley generators.

We do not assert the general-\(n\) isomorphism.  The open general-\(n\) ingredients are listed in Section~\ref{sec:open-points}.

\section{Conventions: arcs, states, diagonal labels, and ground rings}
\label{sec:conventions}
% RECORDS: 397, 699, 1477

Let \(h\) be an invertible formal parameter and put
\[
q=h^2.
\]
The integral ground ring is
\[
\Ring,
\]
and the normalized Chevalley form will use the field
\[
\Field.
\]
Over \(\Ring\), the element
\[
q-q^{-1}=h^2-h^{-2}
\]
is not a unit; over \(\Field\) it is invertible.

For \(n=2\) we use the state set
\[
\{1,2\}=\{+,-\},
\qquad
1=+,
\quad
2=-.
\]
The once-punctured bigon \(\bigon\) is drawn with horizontal arcs oriented to the right.  The lower arcs are denoted \(b_{ij}\), and the upper arcs are denoted \(c_{kl}\).  The first index is the source state and the second index is the target state.
The reduced stated-skein algebra imposes the two bad-arc relations
\[
b_{21}=0,
\qquad
c_{12}=0.
\]
The surviving lower arcs form the positive Borel half, and the surviving upper arcs form the negative Borel half.  In Korinman's notation, the lower arcs are the \(\alpha\)-entries and the upper arcs are the \(\beta\)-entries.  The rank-two dictionary is
\[
\begin{aligned}
b_{11}&=\alpha_{++}, & b_{12}&=\alpha_{+-}, & b_{21}&=\alpha_{-+}=0, & b_{22}&=\alpha_{--},\\
c_{11}&=\beta_{++}, & c_{12}&=\beta_{+-}=0, & c_{21}&=\beta_{-+}, & c_{22}&=\beta_{--}.
\end{aligned}
\]
Set
\[
x=b_{22}=\alpha_{--},
\qquad
y=c_{22}=\beta_{--},
\qquad
u=b_{12}=\alpha_{+-},
\qquad
v=c_{21}=\beta_{-+}.
\]
In the reduced presentation proved in Proposition~\ref{prop:reduced-presentation}, the determinant relations give
\[
b_{11}=x^{-1},
\qquad
c_{11}=y^{-1}.
\]
Thus the six surviving arc generators are
\[
x,\ x^{-1},\ u
\quad\text{on the lower, positive Borel side,}
\]
and
\[
y,\ y^{-1},\ v
\quad\text{on the upper, negative Borel side.}
\]

\section{The reduced rank‑two stated‑skein presentation}
\label{sec:reduced-skein}
% RECORDS: 627, 702, 699, 1477

Korinman's unreduced presentation is generated by the entries of the two \(2\times 2\) matrices
\[
M(\alpha)=
\begin{pmatrix}
\alpha_{++} & \alpha_{+-}\\
\alpha_{-+} & \alpha_{--}
\end{pmatrix},
\qquad
M(\beta)=
\begin{pmatrix}
\beta_{++} & \beta_{+-}\\
\beta_{-+} & \beta_{--}
\end{pmatrix}.
\]
The master relation is Korinman's Equation~(3) \cite{Korinman}:
\[
\mathscr R\bigl(M(Y)\odot M(X)\bigr)
=
\bigl(M(X)\odot M(Y)\bigr)\mathscr R,
\]
for
\[
(X,Y)\in\{(\alpha,\alpha),(\beta,\beta),(\alpha,\beta)\},
\]
together with the quantum determinant conditions
\[
\det_q M(\alpha)=\det_q M(\beta)=1,
\qquad
\det_q
\begin{pmatrix}
a&b\\ c&d
\end{pmatrix}
=
ad-q^{-1}bc.
\]
Here \(\odot\) denotes the Kronecker product of matrices.  Corollary~2.15 of \cite{Korinman} is obtained by expanding this matrix relation entry by entry.  The \((\alpha,\alpha)\) and \((\beta,\beta)\) cases give the same-family relations; the \((\alpha,\beta)\) case gives the mixed exchange law.

% RECORDS: 627, 702, 699, 1477
\begin{proposition}\label{prop:reduced-presentation}
Over \(\Ring\), the reduced stated-skein algebra \(\hS_2(\bigon)\) is generated by
\[
x=b_{22},\quad y=c_{22},\quad u=b_{12},\quad v=c_{21},
\]
with \(x,y\) invertible, with
\[
b_{21}=0,\quad c_{12}=0,\quad b_{11}=x^{-1},\quad c_{11}=y^{-1},
\]
and with defining relations
\begin{align}
xy&=yx, \label{eq:xy-comm}\\
u x&=q^{-1}x u, \label{eq:ux}\\
u y&=q y u, \label{eq:uy}\\
v x&=q x v, \label{eq:vx}\\
y v&=q v y, \label{eq:yv}\\
v u-u v&=-(q-q^{-1})(x^{-1}y-y^{-1}x). \label{eq:vu-uv}
\end{align}
\end{proposition}

\begin{proof}
This is the reduced form of Korinman's presentation of the stated-skein algebra of the once-punctured bigon \cite{Korinman}.  The unreduced relations are those of Corollary~2.15, which are the entrywise expansion of Equation~(3).  The hypotheses of that presentation are the once-punctured bigon, the integral ground ring \(\mathbb Z[A^{\pm1}]\), and the parameter convention \(q=A^2\); in our notation \(A=h\).  The reduced algebra is obtained by adding the bad-arc relations
\[
\alpha_{-+}=0,
\qquad
\beta_{+-}=0.
\]
Under the arc dictionary, these are exactly
\[
b_{21}=0,
\qquad
c_{12}=0.
\]

First consider the determinant relations.  The quantum determinant conditions give
\[
\alpha_{++}\alpha_{--}-q^{-1}\alpha_{+-}\alpha_{-+}=1,
\qquad
\beta_{++}\beta_{--}-q^{-1}\beta_{+-}\beta_{-+}=1.
\]
After imposing \(\alpha_{-+}=0\) and \(\beta_{+-}=0\), these become
\[
\alpha_{++}\alpha_{--}=1,
\qquad
\beta_{++}\beta_{--}=1.
\]
To see that these inverses are two-sided, use the same-family diagonal commutator relations contained in Corollary~2.15, equivalently obtained by expanding Equation~(3):
\[
\alpha_{--}\alpha_{++}-\alpha_{++}\alpha_{--}
=
(q-q^{-1})\alpha_{+-}\alpha_{-+},
\]
and
\[
\beta_{--}\beta_{++}-\beta_{++}\beta_{--}
=
(q-q^{-1})\beta_{+-}\beta_{-+}.
\]
After reduction the right-hand sides vanish, so \(\alpha_{++}\) commutes with \(\alpha_{--}\), and \(\beta_{++}\) commutes with \(\beta_{--}\).  Hence
\[
\alpha_{--}\alpha_{++}=1,
\qquad
\beta_{--}\beta_{++}=1.
\]
Thus
\[
b_{11}=\alpha_{++}=x^{-1},
\qquad
c_{11}=\beta_{++}=y^{-1},
\]
and \(x\) and \(y\) are two-sided units.
The same-family relations also give the reduced same-family \(q\)-commutations
\[
u x=q^{-1}x u,
\qquad
y v=q v y.
\]
The mixed relations give the torus-torus commutation
\[
xy=yx,
\]
and the two surviving mixed off-diagonal \(q\)-commutations
\[
u y=q y u,
\qquad
v x=q x v.
\]
In the notation of Corollary~2.15, the relations in the killed family become \(0=0\) after reduction; the surviving mixed \(q\)-commutations arise from those mixed relations whose right-hand sides contain a killed generator and therefore vanish.  The surviving mixed commutator relation is
\[
v u-u v=-(q-q^{-1})(x^{-1}y-y^{-1}x).
\]
The remaining same-family \(x\)-block relations (for instance \(u x^{-1}=q x^{-1}u\) and \(v y^{-1}=q y^{-1}v\)) are equivalent to the displayed ones under the invertibility of \(x,y\).  Since Corollary~2.15 gives a complete presentation of the unreduced algebra, adding the bad-arc relations and simplifying as above gives precisely the displayed presentation.
\end{proof}

\section{The generalized double and torus calculus}
\label{sec:double-torus}

\subsection{Group-like tori commute in a Majid skew double}
% RECORDS: 913
\begin{proposition}\label{prop:grouplikes-commute}
Let \(H_-\) and \(H_+\) be Hopf algebras over a field, equipped with a skew Hopf pairing \(\tau:H_-\otimes H_+\to \mathbb{k}\).  In the associated generalized double \(H_-\bipr H_+\), any group-like element \(x\in H_-\) and any group-like element \(y\in H_+\) commute.
\end{proposition}

\begin{proof}
In Majid's construction of the double cross product associated to a skew pairing \cite{Majid}, the cross relation specialized to group-like elements has the form
\[
xy=\tau(x,y)\,\tau(Sx,y)\,yx.
\]
Since \(x\) is group-like, \(Sx=x^{-1}\).  The Hopf-pairing axioms give, because \(y\) is group-like,
\[
\tau(xx^{-1},y)=\tau(x,y)\tau(x^{-1},y)=\tau(1,y)=1.
\]
Thus
\[
\tau(x^{-1},y)=\tau(x,y)^{-1}.
\]
Therefore
\[
\tau(x,y)\tau(Sx,y)
=
\tau(x,y)\tau(x^{-1},y)
=
\tau(x,y)\tau(x,y)^{-1}
=
1.
\]
Hence \(xy=yx\).
\end{proof}

\begin{remark}
Thus the nontrivial values of the Drinfeld--Tanisaki pairing on torus generators do not make the two half-Cartan tori noncommute.  Those values appear instead in torus--root-vector cross relations.
\end{remark}

\subsection{The rank‑two generalized double}
% RECORDS: 41, 627, 913, 1477
\begin{proposition}\label{prop:rank-two-double}
Let \(\DqB\) be the \(\Ring\)-algebra generated by
\[
\Khalf,\ \Kmh,\ \Lhalf,\ \Lmh,\ \Ebar,\ \Fbar
\]
subject to the relations
\begin{align}
&\Khalf\Kmh=\Kmh\Khalf=1,
\qquad
\Lhalf\Lmh=\Lmh\Lhalf=1, \label{eq:D-torus-inv}\\
&\text{all elements of }\{\Khalf,\Kmh,\Lhalf,\Lmh\}\text{ commute pairwise}, \label{eq:D-torus-comm}\\
&\Ebar\Khalf=q^{-1}\Khalf\Ebar,
\qquad
\Ebar\Lhalf=q\Lhalf\Ebar, \label{eq:D-E-twist}\\
&\Fbar\Khalf=q\Khalf\Fbar,
\qquad
\Fbar\Lhalf=q^{-1}\Lhalf\Fbar, \label{eq:D-F-twist}\\
&\Ebar\Fbar-\Fbar\Ebar=(q-q^{-1})(K-L), \label{eq:D-comm}
\end{align}
where
\[
K=(\Khalf)^2,
\qquad
L=(\Lhalf)^2.
\]
After scalar extension to \(\Field\), this algebra is the rank-two generalized Drinfeld--Tanisaki double
\[
\check U_q(\mathfrak b_-)\bipr \check U_q(\mathfrak b_+)
\]
with respect to the rank-two specialization of the problem's pairing.  Over \(\Ring\) it is the corresponding integral barred form.
\end{proposition}

\begin{proof}
The algebra with the displayed presentation is precisely the algebra \(\DqB\) of Korinman, Definition~2.18 \cite{Korinman}, with Korinman's parameter \(A\) specialized to our \(h\), so that \(q=A^2=h^2\).  Korinman defines \(\DqB\) as the Drinfeld double of the quantum Borel algebra in rank two, and the presentation above is the integral barred form of that double.

The rank-two specialization of the problem's pairing is
\[
\langle \widehat K_i,\widehat K_j\rangle
=
q^{1/2-\delta_{ij}},
\]
so the torus-pairing matrix is
\[
\begin{pmatrix}
q^{-1/2} & q^{1/2}\\
q^{1/2} & q^{-1/2}
\end{pmatrix}
=
\begin{pmatrix}
h^{-1} & h\\
h & h^{-1}
\end{pmatrix},
\]
with determinant
\[
q^{-1}-q.
\]
Also
\[
\langle E_1,F_1\rangle
=
\frac{1}{q^{-1}-q}
=
-\frac{1}{q-q^{-1}}.
\]
This sign is the standard convention and is consistent with \eqref{eq:D-comm}.
The generalized double of the two rank-two Borel Hopf algebras with this skew pairing is given by Majid's double-cross product construction \cite{Majid}.  Substituting the above pairing values into the cross relations gives the torus--root-vector twists \eqref{eq:D-E-twist}--\eqref{eq:D-F-twist} and the commutator relation \eqref{eq:D-comm}; this is the presentation appearing in Korinman's Definition~2.18.  Proposition~\ref{prop:grouplikes-commute} gives the pairwise commutation of the group-like torus generators \eqref{eq:D-torus-comm}.  Over \(\Field\) the determinant \(q^{-1}-q\) is nonzero, so the usual nondegenerate specialization of the pairing is recovered.  Over \(\Ring\) the same presentation defines the integral barred form.
\end{proof}

The positive Borel half is generated by \(\Khalf\) and \(\Ebar\), while the negative Borel half is generated by \(\Lhalf\) and \(\Fbar\).  In particular, \(\DqB\) is generated by its two Borel halves.

After extending scalars to \(\Field\), one may normalize the Chevalley generators by
\[
E=\frac{\Ebar}{q-q^{-1}},
\qquad
F=\frac{\Fbar}{q-q^{-1}}.
\]
Then \eqref{eq:D-comm} becomes
\[
EF-FE=\frac{K-L}{q-q^{-1}}.
\]
Over \(\Ring\) only the barred form is available.

\subsection{A torus-centrality criterion}
\label{sec:centrality-criterion}
% RECORDS: 1033
The following proposition is purely algebraic and applies to the proposed rank-\(n\) double once the relevant generation hypothesis is satisfied.  It gives the centrality criterion for anti-diagonal torus products.

Let
\[
\Lambda=\mathbb Z^n,
\qquad
c=(1,\dots,1)\in\Lambda,
\]
and define the symmetric bilinear form
\[
B(x,y)
=
\sum_{i=1}^n x_i y_i
-
\frac{\bigl(\sum_{i=1}^n x_i\bigr)\bigl(\sum_{i=1}^n y_i\bigr)}{n}.
\]
Let
\[
\Lambda_0=\bigl\{a\in\Lambda:\sum_i a_i=0\bigr\}
\]
be the root lattice.  If \(a\in\Lambda_0\), then
\[
B(x,a)=\sum_i x_i a_i\in\mathbb Z.
\]

\begin{proposition}\label{prop:centrality-criterion}
Let \(\mathbb{k}\) be a field of characteristic \(0\), and let \(q\in\mathbb{k}^\times\) not be a root of unity.  Let \(W\subset\Lambda_0\) be a set of weights spanning \(\Lambda_0\).  Suppose \(\mathcal A\) is a \(\mathbb{k}\)-algebra generated by commuting invertible torus elements
\[
T^+_\lambda,\ T^-_\lambda
\qquad
(\lambda\in\Lambda)
\]
and root-vector generators
\[
X^+_\alpha,\ X^-_\alpha
\qquad
(\alpha\in W),
\]
with at least one of \(X^+_\alpha,X^-_\alpha\) nonzero for each \(\alpha\in W\).  Suppose the torus-root relations are
\[
T^\pm_\lambda X^+_\alpha (T^\pm_\lambda)^{-1}
=
q^{B(\lambda,\alpha)}X^+_\alpha,
\qquad
T^\pm_\lambda X^-_\alpha (T^\pm_\lambda)^{-1}
=
q^{-B(\lambda,\alpha)}X^-_\alpha.
\]
Then for \(\mu,\nu\in\Lambda\), the element
\[
z(\mu,\nu)=T^-_\mu T^+_\nu
\]
is central in \(\mathcal A\) if and only if
\[
\mu+\nu\in \mathbb Z c.
\]
\end{proposition}

\begin{proof}
Because \(\mathcal A\) is generated by the torus elements and the root-vector generators, it is enough to check commutation with these generators.  Since the torus elements commute with each other, \(z(\mu,\nu)\) commutes with every torus generator.

Let \(\alpha\in W\).  Using the torus-root relations,
\[
\begin{aligned}
z(\mu,\nu)X^+_\alpha z(\mu,\nu)^{-1}
&=
T^-_\mu T^+_\nu X^+_\alpha (T^+_\nu)^{-1}(T^-_\mu)^{-1}\\
&=
T^-_\mu\bigl(q^{B(\nu,\alpha)}X^+_\alpha\bigr)(T^-_\mu)^{-1}\\
&=
q^{B(\nu,\alpha)}q^{B(\mu,\alpha)}X^+_\alpha\\
&=
q^{B(\mu+\nu,\alpha)}X^+_\alpha.
\end{aligned}
\]
Similarly,
\[
z(\mu,\nu)X^-_\alpha z(\mu,\nu)^{-1}
=
q^{-B(\mu+\nu,\alpha)}X^-_\alpha.
\]

If \(\mu+\nu\in\mathbb Z c\), then for every \(a\in\Lambda_0\),
\[
B(\mu+\nu,a)=0,
\]
because \(B(c,a)=\sum_i a_i=0\).  Hence \(z(\mu,\nu)\) commutes with every root-vector generator and therefore with all of \(\mathcal A\).

Conversely, suppose \(z(\mu,\nu)\) is central.  Then for each \(\alpha\in W\),
\[
\bigl(q^{B(\mu+\nu,\alpha)}-1\bigr)X^+_\alpha=0,
\qquad
\bigl(q^{-B(\mu+\nu,\alpha)}-1\bigr)X^-_\alpha=0.
\]
By hypothesis, at least one of \(X^+_\alpha,X^-_\alpha\) is nonzero.  Since \(\mathbb{k}\) is a field and \(q\) is not a root of unity, this forces
\[
B(\mu+\nu,\alpha)=0
\qquad
\text{for every }\alpha\in W.
\]
Because \(W\) spans \(\Lambda_0\), we get
\[
B(\mu+\nu,a)=0
\qquad
\text{for every }a\in\Lambda_0.
\]
For \(a=e_i-e_j\in\Lambda_0\),
\[
B(\mu+\nu,a)=(\mu_i+\nu_i)-(\mu_j+\nu_j)=0
\]
for all \(i,j\).  Thus all coordinates of \(\mu+\nu\) are equal, so
\[
\mu+\nu\in\mathbb Z c.
\]
This proves the equivalence.
\end{proof}

\begin{remark}
The generalized double satisfies the generation hypothesis in the proposition because it is generated by its two Borel halves.  In the rank-two case, however, we next prove the needed centrality statement directly from the defining relations of \(\DqB\).
\end{remark}

% RECORDS: 41, 1033
\begin{corollary}\label{cor:KL-central}
In \(\DqB\), the element
\[
\Khalf\Lhalf
\]
is central.
\end{corollary}

\begin{proof}
It is immediate from \eqref{eq:D-torus-comm} that \(\Khalf\Lhalf\) commutes with all torus generators.  We check the root generators.

From \eqref{eq:D-E-twist}, we have
\[
\Ebar\Khalf=q^{-1}\Khalf\Ebar,
\qquad
\Ebar\Lhalf=q\Lhalf\Ebar.
\]
Equivalently,
\[
\Khalf\Ebar=q\Ebar\Khalf,
\qquad
\Lhalf\Ebar=q^{-1}\Ebar\Lhalf.
\]
Therefore
\[
\begin{aligned}
(\Khalf\Lhalf)\Ebar
&=
\Khalf(\Lhalf\Ebar)\\
&=
\Khalf(q^{-1}\Ebar\Lhalf)\\
&=
q^{-1}(\Khalf\Ebar)\Lhalf\\
&=
q^{-1}(q\Ebar\Khalf)\Lhalf\\
&=
\Ebar(\Khalf\Lhalf).
\end{aligned}
\]

From \eqref{eq:D-F-twist}, we have
\[
\Fbar\Khalf=q\Khalf\Fbar,
\qquad
\Fbar\Lhalf=q^{-1}\Lhalf\Fbar.
\]
Equivalently,
\[
\Khalf\Fbar=q^{-1}\Fbar\Khalf,
\qquad
\Lhalf\Fbar=q\Fbar\Lhalf.
\]
Therefore
\[
\begin{aligned}
(\Khalf\Lhalf)\Fbar
&=
\Khalf(\Lhalf\Fbar)\\
&=
\Khalf(q\Fbar\Lhalf)\\
&=
q(\Khalf\Fbar)\Lhalf\\
&=
q(q^{-1}\Fbar\Khalf)\Lhalf\\
&=
\Fbar(\Khalf\Lhalf).
\end{aligned}
\]
Thus \(\Khalf\Lhalf\) commutes with all generators of \(\DqB\), and hence is central.
\end{proof}

\section{Explicit rank‑two isomorphism}
\label{sec:isomorphism}
% RECORDS: 41, 627, 699, 1477

We now construct the explicit map from the rank-two generalized double to the reduced stated-skein algebra.

Define an assignment on the generators of \(\DqB\) by
\begin{equation}\label{eq:Psi-def}
\begin{aligned}
\Khalf &\longmapsto b_{22}=x,\\
\Kmh &\longmapsto b_{11}=x^{-1},\\
\Lhalf &\longmapsto c_{22}=y,\\
\Lmh &\longmapsto c_{11}=y^{-1},\\
\Ebar &\longmapsto -h\,b_{12}b_{22}=-h\,u x,\\
\Fbar &\longmapsto h^{-1}c_{22}c_{21}=h^{-1}y v.
\end{aligned}
\end{equation}

% RECORDS: 41, 627, 1477
\begin{lemma}\label{lem:psi-respects}
The assignment \eqref{eq:Psi-def} satisfies the defining relations of \(\DqB\) inside \(\hS_2(\bigon)\).
\end{lemma}

\begin{proof}
The torus relations are immediate from
\[
x x^{-1}=x^{-1}x=1,
\qquad
y y^{-1}=y^{-1}y=1,
\qquad
xy=yx.
\]

We check the four torus--root-vector relations.  First, using \(u x=q^{-1}x u\), equivalently \(x u=q u x\),
\[
\Psi(\Ebar)\Psi(\Khalf)
=
(-h u x)x
=
-h u x^2,
\]
while
\[
q^{-1}\Psi(\Khalf)\Psi(\Ebar)
=
q^{-1}x(-h u x)
=
-h q^{-1}x u x
=
-h q^{-1}(q u x)x
=
-h u x^2.
\]
Thus
\[
\Psi(\Ebar)\Psi(\Khalf)
=
q^{-1}\Psi(\Khalf)\Psi(\Ebar).
\]

Second, using \(u y=q y u\) and \(xy=yx\),
\[
\Psi(\Ebar)\Psi(\Lhalf)
=
(-h u x)y
=
-h u x y
=
-h u y x
=
-h(q y u)x
=
q y(-h u x),
\]
so
\[
\Psi(\Ebar)\Psi(\Lhalf)
=
q\Psi(\Lhalf)\Psi(\Ebar).
\]

Third, using \(v x=q x v\) and \(xy=yx\),
\[
\Psi(\Fbar)\Psi(\Khalf)
=
(h^{-1}y v)x
=
h^{-1}y v x
=
h^{-1}y(q x v)
=
q x(h^{-1}y v),
\]
so
\[
\Psi(\Fbar)\Psi(\Khalf)
=
q\Psi(\Khalf)\Psi(\Fbar).
\]

Fourth, using \(y v=q v y\), equivalently \(v y=q^{-1}y v\),
\[
\Psi(\Fbar)\Psi(\Lhalf)
=
(h^{-1}y v)y
=
h^{-1}y v y
=
h^{-1}y(q^{-1}y v)
=
q^{-1}y(h^{-1}y v),
\]
so
\[
\Psi(\Fbar)\Psi(\Lhalf)
=
q^{-1}\Psi(\Lhalf)\Psi(\Fbar).
\]

It remains to check the commutator relation.  We compute
\[
\Psi(\Ebar)\Psi(\Fbar)
=
(-h u x)(h^{-1}y v)
=
-u x y v,
\]
and
\[
\Psi(\Fbar)\Psi(\Ebar)
=
(h^{-1}y v)(-h u x)
=
-y v u x.
\]
Therefore
\[
\Psi(\Ebar)\Psi(\Fbar)-\Psi(\Fbar)\Psi(\Ebar)
=
-u x y v+y v u x.
\]
Using \(xy=yx\), \(u y=q y u\), and \(u x=q^{-1}x u\),
\[
u x y v
=
u y x v
=
(u y)x v
=
q y u x v
=
q y(q^{-1}x u)v
=
y x u v.
\]
Similarly, using \(u x=q^{-1}x u\) and \(v x=q x v\),
\[
y v u x
=
y v(u x)
=
y v(q^{-1}x u)
=
q^{-1}y(v x)u
=
q^{-1}y(q x v)u
=
y x v u.
\]
Thus
\[
\Psi(\Ebar)\Psi(\Fbar)-\Psi(\Fbar)\Psi(\Ebar)
=
y x(v u-u v).
\]
By \eqref{eq:vu-uv},
\[
v u-u v=-(q-q^{-1})(x^{-1}y-y^{-1}x).
\]
Hence
\[
\begin{aligned}
y x(v u-u v)
&=
-(q-q^{-1})y x(x^{-1}y-y^{-1}x)\\
&=
-(q-q^{-1})(y^2-x^2)\\
&=
(q-q^{-1})(x^2-y^2).
\end{aligned}
\]
Under the assignment,
\[
x^2=\Psi(K),
\qquad
y^2=\Psi(L).
\]
Therefore
\[
\Psi(\Ebar)\Psi(\Fbar)-\Psi(\Fbar)\Psi(\Ebar)
=
(q-q^{-1})(\Psi(K)-\Psi(L)),
\]
which is precisely \eqref{eq:D-comm}.
\end{proof}

% RECORDS: 41, 627
\begin{proposition}\label{prop:psi-hom}
The assignment \eqref{eq:Psi-def} extends uniquely to a \(\Ring\)-algebra homomorphism
\[
\Psi:\DqB\longrightarrow \hS_2(\bigon).
\]
This homomorphism is surjective.
\end{proposition}

\begin{proof}
Existence and uniqueness follow from Lemma~\ref{lem:psi-respects} by the universal property of the presentation of \(\DqB\).

For surjectivity, the image contains
\[
x=\Psi(\Khalf),
\qquad
y=\Psi(\Lhalf).
\]
Moreover,
\[
u=-h^{-1}\Psi(\Ebar)x^{-1},
\]
because
\[
\Psi(\Ebar)=-h u x.
\]
Similarly,
\[
v=h y^{-1}\Psi(\Fbar),
\]
because
\[
\Psi(\Fbar)=h^{-1}y v.
\]
Thus the image contains \(x,y,u,v\), and hence all generators of \(\hS_2(\bigon)\) from Proposition~\ref{prop:reduced-presentation}.  Therefore \(\Psi\) is surjective.
\end{proof}

% RECORDS: 41, 627, 699, 1477
\begin{theorem}\label{thm:rank-two-double-skein}
The homomorphism
\[
\Psi:\DqB\longrightarrow \hS_2(\bigon)
\]
is an algebra isomorphism.  With the standard Hopf structures used by Korinman, it is a Hopf isomorphism.
\end{theorem}

\begin{proof}
This is Theorem~2.20 of Korinman \cite{Korinman}, also appearing as Theorem~1.2 in the introduction, translated into the present notation.  Korinman works with a parameter \(A\) satisfying \(q=A^2\); here \(h=A\).  His reduced stated-skein algebra of the once-punctured bigon is obtained from the unreduced one by adding the bad-arc relations
\[
\alpha_{-+}=0,
\qquad
\beta_{+-}=0,
\]
which are exactly
\[
b_{21}=0,
\qquad
c_{12}=0.
\]
His algebra \(\DqB\) has the presentation \eqref{eq:D-torus-inv}--\eqref{eq:D-comm}.  His isomorphism is given by
\[
\Khalf\mapsto \alpha_{--},
\qquad
\Kmh\mapsto \alpha_{++},
\qquad
\Lhalf\mapsto \beta_{--},
\qquad
\Lmh\mapsto \beta_{++},
\]
\[
\Ebar\mapsto -A\,\alpha_{+-}\alpha_{--},
\qquad
\Fbar\mapsto A^{-1}\beta_{--}\beta_{-+}.
\]
Under the dictionary
\[
\alpha_{--}=b_{22}=x,
\quad
\alpha_{++}=b_{11}=x^{-1},
\quad
\alpha_{+-}=b_{12}=u,
\]
\[
\beta_{--}=c_{22}=y,
\quad
\beta_{++}=c_{11}=y^{-1},
\quad
\beta_{-+}=c_{21}=v,
\]
this is exactly the map \(\Psi\) above.
The hypotheses of Korinman's theorem are satisfied: the ground ring is \(\Ring=\mathbb Z[h^{\pm1}]\), the parameter convention is \(q=h^2\), the algebra on the skein side is the reduced once-punctured bigon, and the algebra on the double side is the rank-two generalized Drinfeld--Tanisaki double in integral barred form.  Lemma~\ref{lem:psi-respects} and Proposition~\ref{prop:psi-hom} verify the same relation matching and surjectivity explicitly in the present notation.  Korinman's theorem supplies injectivity and Hopf compatibility.  Hence \(\Psi\) is a Hopf algebra isomorphism.
\end{proof}

After scalar extension to \(\Field\), the normalized Chevalley generators satisfy
\[
E=\frac{\Ebar}{q-q^{-1}},
\qquad
F=\frac{\Fbar}{q-q^{-1}},
\]
and their images are
\[
\Psi(E)= -\frac{h}{q-q^{-1}}\,u x,
\qquad
\Psi(F)= \frac{h^{-1}}{q-q^{-1}}\,y v.
\]

\section{The simply connected quotient at rank two}
\label{sec:quotient}
% RECORDS: 41, 397, 1477, 699

Let
\[
I=(b_{11}c_{11}-1,\ b_{22}c_{22}-1)
\]
be the two-sided ideal in \(\hS_2(\bigon)\) generated by the rank-two instances of the problem's elements \(b_{ii}c_{ii}-1\).  Let
\[
J=(\Khalf\Lhalf-1)
\]
be the corresponding ideal in \(\DqB\).

% RECORDS: 1477, 699
\begin{lemma}\label{lem:ideal-principal}
In \(\hS_2(\bigon)\),
\[
b_{22}c_{22}=xy
\]
is central, and
\[
I=(xy-1).
\]
Equivalently, the two generators \(b_{11}c_{11}-1\) and \(b_{22}c_{22}-1\) generate the same ideal up to a unit.
\end{lemma}

\begin{proof}
First we show that \(xy\) is central.  It commutes with \(x\) and \(y\) because \(xy=yx\).  For \(u\),
\[
u x y
=
(u x)y
=
q^{-1}x u y
=
q^{-1}x(q y u)
=
x y u.
\]
For \(v\),
\[
x y v
=
x(y v)
=
x(q v y)
=
q x v y.
\]
From \(v x=q x v\), equivalently \(x v=q^{-1}v x\), we get
\[
q x v y
=
q(q^{-1}v x)y
=
v x y.
\]
Thus \(xy\) commutes with all generators of \(\hS_2(\bigon)\).

Now
\[
b_{22}c_{22}=xy,
\]
and
\[
b_{11}c_{11}=x^{-1}y^{-1}=(xy)^{-1}.
\]
Therefore
\[
b_{11}c_{11}-1
=
(xy)^{-1}-1
=
-(xy)^{-1}(xy-1).
\]
Since \((xy)^{-1}\) is a unit, the ideals generated by \(b_{11}c_{11}-1\) and \(b_{22}c_{22}-1\) coincide with the principal ideal \((xy-1)\).
\end{proof}

We now define the simply connected rank-two quantum group explicitly.  Let \(\Uqsc\) be the \(\Field\)-algebra generated by
\[
t^{\pm1},\ \mathsf E,\ \mathsf F
\]
subject to the relations
\[
t\mathsf E t^{-1}=q\mathsf E,
\qquad
t\mathsf F t^{-1}=q^{-1}\mathsf F,
\qquad
\mathsf E\mathsf F-\mathsf F\mathsf E
=
\frac{t^2-t^{-2}}{q-q^{-1}}.
\]
This is the usual simply connected rank-two form; see for example Klimyk--Schmüdgen \cite{KlimykSchmudgen}.

% RECORDS: 41, 1477
\begin{lemma}\label{lem:quotient-is-Uqsc}
After extending scalars to \(\Field\),
\[
\bigl(\DqB\otimes_{\Ring}\Field\bigr)\,/
(\Khalf\Lhalf-1)
\cong
\Uqsc.
\]
\end{lemma}

\begin{proof}
Let
\[
Q=\bigl(\DqB\otimes_{\Ring}\Field\bigr)\,/
(\Khalf\Lhalf-1).
\]
In \(Q\), the relation \(\Khalf\Lhalf=1\) gives \(\Lhalf=\Kmh\).  Put
\[
\tilde t=\text{class of }\Khalf.
\]
Then the class of \(\Lhalf\) is \(\tilde t^{-1}\), and
\[
K=\tilde t^2,
\qquad
L=\tilde t^{-2}.
\]
Define elements of \(Q\) by
\[
\mathsf e=\frac{\text{class of }\Ebar}{q-q^{-1}},
\qquad
\mathsf f=\frac{\text{class of }\Fbar}{q-q^{-1}}.
\]
This makes sense because \(q-q^{-1}\) is invertible in \(\Field\).

From \(\Ebar\Khalf=q^{-1}\Khalf\Ebar\), we get
\[
\tilde t\,\mathsf e\,\tilde t^{-1}=q\mathsf e.
\]
From \(\Fbar\Khalf=q\Khalf\Fbar\), we get
\[
\tilde t\,\mathsf f\,\tilde t^{-1}=q^{-1}\mathsf f.
\]
The commutator relation gives
\[
\Ebar\Fbar-\Fbar\Ebar
=
(q-q^{-1})(\tilde t^2-\tilde t^{-2}),
\]
and hence
\[
\mathsf e\mathsf f-\mathsf f\mathsf e
=
\frac{\tilde t^2-\tilde t^{-2}}{q-q^{-1}}.
\]
Thus the assignment
\[
t\mapsto \tilde t,
\qquad
\mathsf E\mapsto \mathsf e,
\qquad
\mathsf F\mapsto \mathsf f
\]
defines a homomorphism
\[
\phi:\Uqsc\longrightarrow Q.
\]

Conversely, define a homomorphism from \(\DqB\otimes_{\Ring}\Field\) to \(\Uqsc\) by
\[
\Khalf\mapsto t,
\qquad
\Kmh\mapsto t^{-1},
\qquad
\Lhalf\mapsto t^{-1},
\qquad
\Lmh\mapsto t,
\]
\[
\Ebar\mapsto (q-q^{-1})\mathsf E,
\qquad
\Fbar\mapsto (q-q^{-1})\mathsf F.
\]
The element \(\Khalf\Lhalf-1\) maps to \(tt^{-1}-1=0\).  The torus relations are immediate.  The relations \(\Ebar\Khalf=q^{-1}\Khalf\Ebar\) and \(\Fbar\Khalf=q\Khalf\Fbar\) follow from
\[
\mathsf E t=q^{-1}t\mathsf E,
\qquad
\mathsf F t=q t\mathsf F.
\]
The relations \(\Ebar\Lhalf=q\Lhalf\Ebar\) and \(\Fbar\Lhalf=q^{-1}\Lhalf\Fbar\) follow from
\[
\mathsf E t^{-1}=q t^{-1}\mathsf E,
\qquad
\mathsf F t^{-1}=q^{-1}t^{-1}\mathsf F.
\]
Finally,
\[
[(q-q^{-1})\mathsf E,(q-q^{-1})\mathsf F]
=
(q-q^{-1})^2\frac{t^2-t^{-2}}{q-q^{-1}}
=
(q-q^{-1})(t^2-t^{-2}),
\]
which is the required image of \eqref{eq:D-comm}.  Hence the map factors through \(Q\), giving a homomorphism
\[
\psi:Q\longrightarrow \Uqsc.
\]
By construction, \(\phi\) and \(\psi\) are inverse on generators.  Therefore \(Q\cong \Uqsc\).
\end{proof}

% RECORDS: 41, 397, 1477
\begin{theorem}\label{thm:rank-two-quotient}
Under the isomorphism \(\Psi\) of Theorem~\ref{thm:rank-two-double-skein}, the ideal
\[
J=(\Khalf\Lhalf-1)
\]
corresponds to the skein ideal
\[
I=(b_{11}c_{11}-1,\ b_{22}c_{22}-1).
\]
Consequently, after extending scalars to \(\Field\),
\[
\bigl(\hS_2(\bigon)\otimes_{\Ring}\Field\bigr)\,/I
\cong
\Uqsc.
\]
Over \(\Ring\), the same quotient gives the barred integral form of the simply connected rank-two quantum group.
\end{theorem}

\begin{proof}
We have
\[
\Psi(\Khalf\Lhalf)
=
\Psi(\Khalf)\Psi(\Lhalf)
=
b_{22}c_{22}
=
xy.
\]
Thus
\[
\Psi(\Khalf\Lhalf-1)=xy-1.
\]
By Lemma~\ref{lem:ideal-principal},
\[
I=(xy-1).
\]
Since \(\Psi\) is an algebra isomorphism, it identifies \(J\) with \(I\) and therefore induces an isomorphism of quotients:
\[
\DqB/J
\cong
\hS_2(\bigon)/I.
\]
After scalar extension to \(\Field\), Lemma~\ref{lem:quotient-is-Uqsc} identifies the double-side quotient with \(\Uqsc\).  Hence
\[
\bigl(\hS_2(\bigon)\otimes_{\Ring}\Field\bigr)\,/I
\cong
\Uqsc.
\]

Over \(\Ring\), where \(q-q^{-1}\) is not inverted, the quotient is the corresponding barred integral form.  On the double side, after imposing \(\Khalf\Lhalf=1\), set \(t=\Khalf\).  Then \(\Lhalf=t^{-1}\), and the quotient is generated over \(\Ring\) by
\[
t^{\pm1},\ \Ebar,\ \Fbar
\]
with relations
\[
t\Ebar t^{-1}=q\Ebar,
\qquad
t\Fbar t^{-1}=q^{-1}\Fbar,
\qquad
\Ebar\Fbar-\Fbar\Ebar=(q-q^{-1})(t^2-t^{-2}).
\]
Transported by \(\Psi\), this is the same as the skein quotient generated by \(x^{\pm1},u,v\), with \(y=x^{-1}\), subject to the reduced skein relations of Proposition~\ref{prop:reduced-presentation}.
\end{proof}

\section{Open points for general \(n\)}
\label{sec:open-points}
% RECORDS: 1033

The general-\(n\) statement of the original problem is not proved here.  The missing ingredients include:
\begin{enumerate}[label=(\arabic*)]
\item the exact general-\(n\) identification of the reduced stated-skein algebra of the bigon with the appropriate generalized double, including body-level verification of the relevant Lê--Sikora bigon identification and matrix coproduct \cite{LeSikora};
\item compatibility of the geometric splitting homomorphism with the reduced quotient;
\item verification that the surviving triangular skein generators satisfy the full Serre relations;
\item proof that the skein ideal generated by the elements \(b_{ii}c_{ii}-1\) corresponds to the intended torus relations for all \(i\);
\item determination of whether the elements \(b_{ii}c_{ii}\) are individually central for \(n\ge 3\);
\item resolution of the ground-ring issues involving \(q^{1/n}\) and \(q^{1/(2n)}\).
\end{enumerate}
The torus-centrality criterion of Proposition~\ref{prop:centrality-criterion} is proved algebraically in the generated-algebra form; its full general-\(n\) skein interpretation is not established in this paper.

\begin{thebibliography}{9}

\bibitem{Korinman}
G.~Korinman,
\emph{Quantum holonomic link invariants derived from stated skein algebras},
arXiv:1907.01732v4.

\bibitem{Majid}
S.~Majid,
\emph{Foundations of Quantum Group Theory},
Cambridge University Press, 1995.

\bibitem{KlimykSchmudgen}
A.~Klimyk and K.~Schmüdgen,
\emph{Quantum Groups and Their Representations},
Springer, 1997.

\bibitem{LeSikora}
T.~T.~Q.~Lê and A.~Sikora,
\emph{Stated \(\mathrm{SL}(n)\)-Skein Modules and Algebras},
arXiv:2201.00045v3.

\end{thebibliography}

\end{document}